\subsection{内部正则性}\label{sec:elliptic-interior-regularity}

在 Section~\ref{sec:dirichlet-existence-and-uniqueness} 中, 已经证明了 $2k$ 阶椭圆算子的 Dirichlet 问题存在弱解 $u\in H^k(\Omega)$. 现在希望证明, 对系数 $a_{\sigma\gamma}$, 强迫函数 $f$ 和 $\Omega$ 的边界的光滑性作适当假设后, 所得弱解实际上是强解或经典解. 为说明准备如何进行, 先作两次形式计算.

第一次计算中, 假设 $\Omega$ 有光滑边界 $\partial\Omega$, 其单位外法向量为 $\bm{\eta}=(\eta_1,\ldots,\eta_n)$, 并且 $u$ 是
\begin{equation}\label{eq:interior-regularity-formal-poisson-equation}
-\Delta u=f
\end{equation}
在 $\Omega$ 中的经典解, 且
\begin{equation}\label{eq:interior-regularity-formal-dirichlet-boundary}
u=0
\end{equation}
在 $\partial\Omega$ 上成立. 目标是证明, 上述椭圆问题的 (弱) 解实际上属于比 $H_0^1(\Omega)$ ``更好\textquotesingle\textquotesingle\ 的空间. 为此作准备, 现在用 $H^1(\Omega)$ 范数估计 $u$ 的二阶偏导数矩阵的 $L^2(\Omega)$ 范数. 因为这只是形式计算, 将假定已经知道 $u$ 具有任意所需的光滑性.

\[\begin{aligned}
\int_\Omega|\Delta u|^2\,\mathrm{d}\bm{x}
&=\sum_{i,j=1}^n\int_\Omega u_{x_ix_i}u_{x_jx_j}\,\mathrm{d}\bm{x}\\
&=-\sum_{i,j=1}^n\int_\Omega u_{x_i}u_{x_jx_ix_j}\,\mathrm{d}\bm{x}\\
&\quad+\sum_{i,j=1}^n\int_{\partial\Omega}u_{x_i}u_{x_jx_j}\eta_i\,\mathrm{d}S\\
&=\sum_{i,j=1}^n\int_\Omega u_{x_ix_j}u_{x_jx_i}\,\mathrm{d}\bm{x}\\
&\quad+\sum_{i,j=1}^n\int_{\partial\Omega}u_{x_i}u_{x_jx_j}\eta_i-u_{x_i}u_{x_ix_j}\eta_j\,\mathrm{d}S.
\end{aligned}\]

还有
\begin{equation}\label{eq:interior-regularity-formal-laplacian-data-identity}
\int_\Omega|\Delta u|^2\,\mathrm{d}\bm{x}=\int_\Omega|f|^2\,\mathrm{d}\bm{x}.
\end{equation}
结合这两个结果得到
\begin{equation}\label{eq:interior-regularity-formal-hessian-boundary-estimate}
\int_\Omega|\nabla^2u|^2\,\mathrm{d}\bm{x}\leq\int_\Omega|f|^2\,\mathrm{d}\bm{x}+|\text{boundary terms}|.
\end{equation}
因此, 若有关于边界项的进一步信息, 就能用数据 $f$ 导出解 $u$ 的 $H^2(\Omega)$ 范数的先验估计.
